% !PRINTED 153
\weberpage{153}
%<<CONT
expressed by the $x$ by means of a new linear substitution $E$, which we may denote thus:

\weq{10}{z = E (x) = B A (x),}

and the substitution $E = BA$ is said to be \gesperrt{composed} or \gesperrt{compounded} of $B$ and $A$. But a distinction must here be made between $AB$ and $BA$. Two substitutions $A, B$ with the special property that $AB = BA$ are said to be \gesperrt{permutable} or \gesperrt{commutative} with one another. The substitution coefficients of $E$ are obtained by inserting the expressions (1) in (9):

% [sic: Weber prints i below the first sum although h is its summation index; 1,n is expected.]
\weq{11}{\begin{gathered} z_k = \sum_{i}^{h} e_h^{(k)} x_h, \\ e_h^{(k)} = \sum_{1,n}^{i} b_i^{(k)} a_h^{(i)}. \end{gathered}}

These formulae are exactly the same as those we used in \S. 27 of the first volume\footnote{Only there, for an easily discernible reason, the notation was somewhat different.}, and we can accordingly express the law of the composition of substitutions expressed in (11) as follows:

\gesperrt{1. To form the substitution $BA$ composed of two substitutions $B, A$, one proceeds exactly as if the two determinants $|B|, |A|$ were to be multiplied together according to the rule of multiplication. Here, in order to form the elements of a row, the elements of a row of the first component are to be multiplied by the corresponding elements of the columns of the second component and then added. This is also expressed briefly by saying that in the first component the summation is taken by rows, in the second by columns.}

From this rule follows the consequence:

\gesperrt{2. The determinant of a composed substitution is equal to the product of the determinants of the components.}
